\documentclass{beamer}

  %Reconhecer os caracteres especiais no linux
  %\usepackage[utf8x]{inputenc}
  %Reconhcer os caracteres especiais no Windows
\usepackage[utf8]{inputenc} %para portugu?s
%  \usepackage[brazil]{babel} %para portugu?s
%\usepackage[latin1]{inputenc} %para portugu?s
%\usepackage[T1]{fontenc} %para portugu?s
  %Recursos matemático como por exemplo símbolos
  \usepackage{amsmath}
  \usepackage{amsfonts}
  \usepackage{amssymb}
  \usepackage{enumerate}
  %Idioma
  %\usepackage[portuges]{babel}
  %Usar gráficos
  \usepackage{graphicx}
  %Usar animação no beamer
  %\usepackage{animate}
  %Para inserir videos
  %\usepackage{movie15}
  %\usepackage{hyperref}

%##############################.##########################################################
% Temas
%########################################################################################
   %    \usepackage{beamerthemeshadow} % Usaremos este
%       \usepackage{beamerthemebars}
%       \usepackage[headheight=12pt,footheight=12pt]{beamerthemeboxes}
%       \usepackage{beamerthemeclassic}
%       \usepackage{beamerthemelined}
%       \usepackage[width=12pt,dark,tab]{beamerthemesidebar}
%       \usepackage{beamerthemesplit}
%       \usepackage{beamerthemetree}
%       \usepackage[bar]{beamerthemetree}
%       \usepackage{beamerthemeWarsaw}
%       \usepackage{beamerthemeMadrid}
%       \usepackage{beamerthemeAnnArbor}
%       \usepackage{beamerthemeAntibes}
%       \usepackage{beamerthemeBergen}
%      \usepackage{beamerthemeBerkeley}
%       \usepackage{beamerthemeBerlin}
%       \usepackage{beamerthemeBoadilla}
%       \usepackage{beamerthemeboxes}
%       \usepackage{beamerthemeCambridgeUS}
%       \usepackage{beamerthemeCopenhagen}
%       \usepackage{beamerthemeDarmstadt}
%       \usepackage{beamerthemedefault}
%       \usepackage{beamerthemeDresden}
       \usepackage{beamerthemeFrankfurt}
%       \usepackage{beamerthemeGoettingen}
%       \usepackage{beamerthemeHannover}
 %     \usepackage{beamerthemeIlmenau}
%       \usepackage{beamerthemeJuanLesPins}
%       \usepackage{beamerthemeLuebeck}
%       \usepackage{beamerthemeMadrid}
%      \usepackage{beamerthemeMalmoe}
%       \usepackage{beamerthemeMarburg}
%       \usepackage{beamerthemeMontpellier}
%       \usepackage{beamerthemePaloAlto}
%       \usepackage{beamerthemePittsburgh} %** seco
%       \usepackage{beamerthemeRochester}
%       \usepackage{beamerthemeSingapore}
%       \usepackage{beamerthemeSzeged}
%      \usepackage{beamerthemeWarsaw}
%       \usepackage{beamerthemeHannover}

%########################################################################################
%Cores que podem ser usadas em diferentes themes
%########################################################################################
      %  \usecolortheme{default}
%	\usecolortheme[RGB={255,0,0}]{structure}
%	\usecolortheme[RGB={0,255,0}]{structure}
% 	\usecolortheme[RGB={0,0,255}]{structure}
%       \usecolortheme{albatross} %azul
    \usecolortheme{beaver} %vermelho
%       \usecolortheme{beetle} %cinza
%       \usecolortheme{crane} %amarelo
%       \usecolortheme{dolphin} %azul claro
%       \usecolortheme{dove} %branco
%       \usecolortheme{fly} %cinza
%       \usecolortheme{lily} % banco
%       \usecolortheme{orchid} % branco
%       \usecolortheme{seagull} %cinza
%       \usecolortheme{seahorse} %azul
%       \usecolortheme{sidebartab} %branco
%       \usecolortheme{whale} %branco
%     \usecolortheme{wolverine} %amarelo e azul
%########################################################################################
%########################################################################################
% Tipo o simbolo na listagem
% \setbeamertemplate{items}[ball]
% \setbeamertemplate{items}[circle]
 \setbeamertemplate{items}[rectangle]
% \setbeamertemplate{items}[default] %triangulos
%########################################################################################
%########################################################################################

%Aqui você digita o título e a mensagem que irá aparecer no rodapé
\title[Mensagem no Rodapé]{Título}
%Aqui você digita o seu nome para aparecer logo abaixo do título
\author{Seu nome}

\begin{document}

  \begin{frame}
    \titlepage %Colocando o título da sua apresentação
  \end{frame}

  % Para o roteiro da apresentação sair antes de cada seção
  \AtBeginSection[]
  {
    \begin{frame}
		\frametitle{Sumário}
		\small
		\tableofcontents[currentsection,hideothersubsections]
		\normalsize
    \end{frame}
  }

  \section{Primeira seção}

  %Ininiando um slide
  \begin{frame}
      \frametitle{Título do Slide} %Colocando o título do slide

      Aqui você entra com o seu texto em \LaTeX!

  \end{frame}

  \section{Segunda seção}

  \begin{frame}
      \frametitle{Um exemplo de animação no Beamer}
      \begin{figure}[t]
	\centering
	\includegraphics<1>[scale=0.5]{1.eps}
	\includegraphics<2>[scale=0.5]{2.eps}
	\includegraphics<3>[scale=0.5]{3.eps}
	\includegraphics<4->[scale=0.5]{4.eps}
      \end{figure}
  \end{frame}

\section{Terceira seção}
  \begin{frame}
    \frametitle{Aparece todos os itens de uma vez}
    \begin{enumerate}%[<+_|alert@+>]
	\item  alerta um
	\item  alerta dois
	\item  alerta tres
	%\pause
	%\item 4  aaaa
	%\item 5  vvvvv
    \end{enumerate}
  \end{frame}

  \begin{frame}
    \frametitle{Aparece um item de cada vez}
    \begin{itemize}
	\item<+-> 1
	\item<+-> 2
	\item<+-> 3
	\item<+-> 4
	\item<+-> 5
    \end{itemize}
  \end{frame}

  \section{Quarta seção}

  \begin{frame}

    \begin{block}{This is a Block}
	This is important information
    \end{block}

    \begin{alertblock}{This is an Alert block}
    This is an important alert
    \end{alertblock}

    \begin{exampleblock}{This is an Example block}
    This is an example
    \end{exampleblock}

  \end{frame}

    \section{Quinta seção}

  \begin{frame}
      \begin{columns}[c] % the "c" option specifies center vertical alignment
      \column{.5\textwidth} % column designated by a command
      Contents of the first column
      \column{.5\textwidth}
      Contents split \\ into two lines
      \end{columns}
  \end{frame}

  \begin{frame}
      \begin{columns}[t] % contents are top vertically aligned
      \begin{column}[T]{5cm} % each column can also be its own environment
      Contents of first column \\ split into two lines
      \end{column}
      \begin{column}[T]{5cm} % alternative top-align that's better for graphics
	    \includegraphics[height=3cm]{1.eps}
      \end{column}
      \end{columns}
  \end{frame}

  %Exemplo de slides limpo
  \frame[plain]{

    Exemplo de slides

  }

  %Exemplo de slide sem cabecalho
  \frame[shrink]{

    Exemplo de slides

  }




\end{document} 